
%%%%%%%%%%%%%%%%%%%%%%%%%%%%%%%%%%%%%%%%%%%%%%%%%%%%%%%%%%%%%%%%%%%%%%%%%%%%%%
% Copyright (c) 2003-2026 by the esys.escript Group
% https://github.com/LutzGross/esys-escript.github.io
%
% Primary Business: Queensland, Australia
% Licensed under the Apache License, version 2.0
% http://www.apache.org/licenses/LICENSE-2.0
%
% See CREDITS file for contributors and development history
%
%%%%%%%%%%%%%%%%%%%%%%%%%%%%%%%%%%%%%%%%%%%%%%%%%%%%%%%%%%%%%%%%%%%%%%%%%%%%%%

\chapter{The \oxley Module}\label{chap:oxley}\index{oxley}

\oxley provides meshes which can be refined locally. A mesh is a forest of
quadtrees in 2D and of octrees in 3D, managed by the \texttt{p4est}
library\index{p4est}: the domain is divided into a grid of rectangular blocks,
and each block is subdivided recursively into quadrants (octants in 3D), more
often where more resolution is needed. The number of subdivisions an element
has gone through is its \emph{level}. Neighbouring elements differ by at most
one level, so where a fine element meets a coarse one a node of the fine
element lies in the middle of an edge (or face) of the coarse element. These
\emph{hanging nodes}\index{hanging node} are what make the mesh
non-conforming.

The \oxley module is used to
\begin{itemize}
\item create a mesh with \function{Rectangle} or \function{Brick},
\item refine it locally with a \class{RefinementQueue2D} or
\class{RefinementQueue3D}, and
\item hand it to \finley to solve PDEs on it, see
Section~\ref{sec:oxley:pdes}.
\end{itemize}
The examples used in this chapter are \file{refinement_queue.py},
\file{refinement_queue_3d.py} and \file{refinement_queue_mpi4py.py} in the
\file{usersguide} directory of the examples.

\section{Meshes}
A rectangle of $4 \times 4$ blocks, each subdivided twice, so that the mesh
has $16 \times 16$ elements, is created by
\begin{python}
from esys.oxley import Rectangle
coarse = Rectangle(n0=4, n1=4, l0=1., l1=1., refine_level=2)
\end{python}
The refinement level can also be given per block, as an array of shape
\var{(n0, n1)}, which grades the mesh across the block seams:
\begin{python}
graded = Rectangle(n0=2, n1=2, l0=1., l1=1., refine_level=[[1, 3], [2, 1]])
\end{python}
\function{Brick} creates a 3D mesh in the same way.

\section{Local Refinement}\label{sec:oxley:refinement}
\index{oxley!refinement}\index{RefinementQueue}
A mesh is refined by collecting refinement operations on a
\class{RefinementQueue2D} (for a \function{Rectangle}) or
\class{RefinementQueue3D} (for a \function{Brick}) and applying them to the
mesh:
\begin{python}
from esys.oxley import RefinementQueue2D
queue = RefinementQueue2D()
queue.setRefinementLevel(4)
queue.refineCircle(x0=0.3, y0=0.6, r=0.1)            # level 4
queue.refineBorder(border="top", dx=0.1, level=3)    # a strip along x_1=1
queue.refineRegion(x0=0.7, y0=0.0, x1=0.9, y1=0.2)   # level 4
fine = queue.apply(coarse)
\end{python}
Three properties of a queue are worth keeping in mind:
\begin{itemize}
\item \method{apply} does not change the domain it is given: it returns a new,
refined domain. The coarse domain, and any \Data object defined on it, stay
valid. \Data objects on the coarse and on the refined domain cannot be
combined; to carry values over, interpolate them or recompute them on the
refined domain.
\item A queue is a template which is not tied to a domain. The same queue can
be applied to any number of domains of the right dimension, and is unchanged
by being applied.
\item A level is absolute: an element is subdivided until it reaches the
requested level, and elements which are already as fine are left alone.
Applying the same queue twice therefore changes nothing; only a higher level
refines further.
\end{itemize}
Each operation takes an optional \var{level}. Without it, the level set by
\method{setRefinementLevel} is used. The operations are applied in the order
in which they were queued; \method{print} lists them.

\subsection{Refinement Driven by Data}\label{sec:oxley:mask}
\index{oxley!mask refinement}
\method{refineMask} refines the elements in which a mask is positive, which
lets the refinement follow a feature of the model rather than a geometric
shape. As a queue is not tied to a domain, it does not hold the mask itself,
only a name for it, the \emph{tag}. The mask is handed to \method{apply} as a
keyword argument named after the tag, and has to be a scalar \Data object on
the domain being refined:
\begin{python}
def sourceMask(domain):
    x = Function(domain).getX()
    return whereNegative(length(x - [0.3, 0.6]) - 0.05)

queue.refineMask("source", level=5)
fine = queue.apply(coarse, source=sourceMask(coarse))
\end{python}
An element is refined if the mask is positive at any of its quadrature points.
To apply the queue to another domain, the mask is created on that domain:
\begin{python}
other = Rectangle(n0=8, n1=8, l0=1., l1=1., refine_level=1)
finer = queue.apply(other, source=sourceMask(other))
\end{python}
\method{apply} raises an exception if a tag in the queue has no mask, if a
keyword argument is not the tag of any mask refinement in the queue (most
likely a misspelt tag), or if a mask is not a scalar on the domain being
refined. A tag must be a valid \PYTHON name other than \var{domain}.

\section{Solving PDEs}\label{sec:oxley:pdes}\index{oxley!solving PDEs}
PDEs on an \oxley mesh are solved by \finley. \method{toFinley} converts the
mesh, hanging nodes included, into a conforming \finley mesh, on which a
\LinearPDE is set up and solved as usual. \function{fromFinleyData} brings
the solution back to the nodes of the \oxley mesh:
\begin{python}
import esys.finley
from esys.oxley import fromFinleyData

def solve(domain):
    fin = domain.toFinley()
    x = fin.getX()
    pde = LinearPDE(fin)
    pde.setSymmetryOn()
    gammaD = whereZero(x[0]) + whereZero(x[0]-1.) \
           + whereZero(x[1]) + whereZero(x[1]-1.)
    pde.setValue(A=kronecker(fin), Y=sourceMask(fin), q=gammaD)
    return fromFinleyData(pde.getSolution(), domain)
\end{python}
Note that \finley has to be imported, even if it is not used by name: without
it \method{toFinley} returns a domain which cannot be used to build a
\LinearPDE. Coefficients can be defined directly on the \finley mesh, as
above, or transferred from the \oxley mesh with \function{toFinleyData}.

A \LinearPDE can also be built directly on an \oxley mesh which has no hanging
nodes, i.e.\ one in which all elements have the same level. On a refined or
graded mesh this raises a \texttt{NotImplementedError}.

\section{Example}\label{sec:oxley:example}
The example \file{refinement_queue.py} solves the Poisson problem
\begin{equation}
-\Delta u = f \mbox{ in } [0,1]^2, \quad u = 0 \mbox{ on the boundary},
\end{equation}
where $f=1$ in a disc of radius $0.05$ around $(0.3,0.6)$ and $f=0$ elsewhere,
once on the coarse mesh of $16 \times 16$ elements and once on a mesh refined
around the source:
\begin{python}
coarse = Rectangle(n0=4, n1=4, l0=1., l1=1., refine_level=2)

queue = RefinementQueue2D()
queue.setRefinementLevel(4)
queue.refineCircle(x0=XC, y0=YC, r=2*R)             # around the source, level 4
queue.refineMask("source", level=5)                 # the source itself, level 5
queue.refineBorder(border="top", dx=0.1, level=3)   # a strip along y=1
queue.refineRegion(x0=0.7, y0=0.0, x1=0.9, y1=0.2)  # a box, the default level
queue.print()

fine = queue.apply(coarse, source=sourceMask(coarse))

u_coarse = solve(coarse)
u_fine = solve(fine)
\end{python}
The output is
\begin{verbatim}
0: Circle at (0.3, 0.6) with r = 0.1, level=4
1: Mask 'source', level=5
2: Boundary  (Top) to depth 0.1
3: Region (0.7, 0) (0.9, 0.2), level=4
coarse mesh: 256 elements, max(u) = 2.510454e-03
refined mesh: 907 elements, max(u) = 3.359582e-03
queue applied to an 8 x 8 block domain: 2899 elements
\end{verbatim}
On the coarse mesh the source is resolved by only a few elements, and its
maximum is underestimated by a quarter. The solutions are written to
\file{refinement_coarse.vtu} and \file{refinement_fine.vtu} in the
subdirectory \file{data}, which the script creates with the \MPI-safe
\function{mkDir}:
\begin{python}
mkDir("data")
saveVTK(os.path.join("data", "refinement_coarse.vtu"), u=u_coarse)
saveVTK(os.path.join("data", "refinement_fine.vtu"), u=u_fine)
\end{python}

\subsection{A 3D Example}
\file{refinement_queue_3d.py} solves the same problem in the unit cube with a
\class{RefinementQueue3D}. The source is a ball of radius $0.1$ around
$(0.3,0.6,0.4)$, refined with \method{refineSphere} and \method{refineMask},
and the layer below the top face, $x_2=1$, is refined with
\method{refineBorder}:
\begin{python}
coarse = Brick(n0=4, n1=4, n2=4, l0=1., l1=1., l2=1., refine_level=1)

queue = RefinementQueue3D()
queue.setRefinementLevel(3)
queue.refineSphere(x0=XC, y0=YC, z0=ZC, r=2*R)      # around the source, level 3
queue.refineMask("source", level=4)                 # the source itself, level 4
queue.refineBorder(border="top", dx=0.1, level=2)   # a layer below x2=1

fine = queue.apply(coarse, source=sourceMask(coarse))
\end{python}
In 3D, \var{"top"} and \var{"bottom"} are the faces normal to $x_2$,
\var{"north"} and \var{"south"} (or \var{"back"} and \var{"front"}) those
normal to $x_1$, and \var{"east"} and \var{"west"} (or \var{"right"} and
\var{"left"}) those normal to $x_0$. The output is
\begin{verbatim}
0: Sphere at (0.3, 0.6, 0.4) with r = 0.2, level=3
1: Mask 'source', level=4
2: Boundary  (Top) to depth 0.1
coarse mesh: 512 elements, max(u) = 3.379138e-03
refined mesh: 6791 elements, max(u) = 4.227389e-03
\end{verbatim}
and the solutions are saved as \file{refinement3d_coarse.vtu} and
\file{refinement3d_fine.vtu} in the subdirectory \file{data}.

\section{Refinement and \MPI}\label{sec:oxley:mpi}\index{oxley!MPI}
A queue has no communicator of its own: \method{apply} copies the domain it is
given, so the refined domain lives on the same communicator as that domain,
and \method{apply} is collective over it. Every rank of the communicator has
to call \method{apply}, with the same queue, which is the case when every rank
runs the same script. A mask is distributed like any other \Data object, each
rank holding the values on its own elements; \method{apply} gathers the marked
elements over the communicator, so that all ranks refine the same ones. The
refined mesh, its \finley export and the solution obtained on it are
therefore the same however many ranks are used.

To refine on a sub-communicator\index{mpi4py!oxley}, it is enough to create the
domain on it, see Chapter~\ref{CHAP:mpi4py}. The example
\file{refinement_queue_mpi4py.py} splits the ranks into up to three groups.
Each group refines the same coarse mesh around the source to a different
level, solves the Poisson problem, and saves its solution as
\file{data/refinement_group<k>.vtu}; world rank 0 collects a report from each
group:
\begin{python}
world = MPI.COMM_WORLD
mkDir("data")      # collective over all ranks, so before the split

ngroups = min(world.size, MAX_GROUPS)
color = world.rank % ngroups
sub = world.Split(color, world.rank)
level = BASE_LEVEL + 1 + color

coarse = Rectangle(n0=4, n1=4, l0=1., l1=1., refine_level=BASE_LEVEL, comm=sub)

queue = RefinementQueue2D()
queue.refineCircle(x0=XC, y0=YC, r=2*R, level=level)
queue.refineMask("source", level=level+1)
fine = queue.apply(coarse, source=sourceMask(coarse))

u = solve(fine)
saveVTK(os.path.join("data", "refinement_group%d.vtu" % color), u=u)
\end{python}
Run on six ranks with
\begin{verbatim}
run-escript -n 6 refinement_queue_mpi4py.py
\end{verbatim}
each group has two ranks, and the output is
\begin{verbatim}
group ranks level  elements        max(u)   integral(u)
    0     2     3       352  3.202207e-03  4.573424e-04
    1     2     4       658  3.357074e-03  4.795395e-04
    2     2     5      1774  3.340626e-03  4.751504e-04
\end{verbatim}
which is the same as with one rank per group. Note that a reduction over all
ranks, such as \function{getMPIWorldSum}, mixes the groups. A sum over the
ranks of a domain uses the communicator the domain lives on, which
\method{getMPIComm} returns. The examples count the elements of a domain this
way, adding up the elements each rank owns, which \method{getMeshInfo}
reports; the samples of a \class{Function} would also count the elements a
rank shares with its neighbours:
\begin{python}
def numElements(domain):
    local = int(domain.getMeshInfo()["numElements"])
    comm = domain.getMPIComm()
    if comm is None:     # without mpi4py the domain lives on all ranks
        return getMPIWorldSum(local)
    return comm.allreduce(local)
\end{python}

\subsection{Node Coordinates at Hanging Nodes}\label{sec:oxley:mastercoords}
\index{oxley!hanging node}\index{hanging node!coordinates}
A hanging node is not a node of the \oxley mesh: its value is the average of
its \emph{masters}, the corners of the coarse edge or face it lies on. A rank
can hold a master without holding any element of which it is a corner, namely
when the coarse element and every fine element meeting at that corner belong
to other ranks; the master then only enters through the hanging nodes it
defines. Its position follows from the fine element alone, since a hanging
node lies half way between a corner of that element and the master, and
\oxley computes it from there. So every rank knows the coordinates of every
node it holds, and \method{getX}, any \Data object built from it, the hanging
node values and the transfers to and from \finley are complete on every rank.

\section{Functions}
\paragraph{Rectangle(n0=10, n1=10, l0=1., l1=1., refine_level=0,
diracPoints=list(), diracTags=list(), comm=None)}
generates a \Domain object representing a two-dimensional rectangle between
$(0,0)$ and $(l0,l1)$, made of \var{n0} by \var{n1} blocks. \var{l0} and
\var{l1} can also be pairs giving the coordinate range along an axis. Each
block is subdivided \var{refine_level} times; \var{refine_level} is either an
integer or an array of shape \var{(n0, n1)} with a level per block.
\var{diracPoints} is a list of coordinate-tuples of points within the mesh,
each point tagged with the respective string within \var{diracTags}.
\var{comm} is an optional \MPI communicator (from \texttt{mpi4py}).

\paragraph{Brick(n0=10, n1=10, n2=10, l0=1., l1=1., l2=1., refine_level=0,
diracPoints=list(), diracTags=list(), comm=None)}
generates a \Domain object representing a three-dimensional brick between
$(0,0,0)$ and $(l0,l1,l2)$, made of \var{n0} by \var{n1} by \var{n2} blocks.
The arguments are as for \function{Rectangle}; a per-block \var{refine_level}
has shape \var{(n0, n1, n2)}.

\paragraph{RefinementQueue2D()} and \textbf{RefinementQueue3D()}
create an empty queue of refinements for a \function{Rectangle} and a
\function{Brick} respectively. The default level of the queue is $0$. A queue
has the following methods; \var{level} defaults to the level of the queue.
\begin{itemize}
\item \method{setRefinementLevel}(\var{level}) sets the default level.
\item \method{refineUniform}(\var{level}) refines every element.
\item \method{refinePoint}(\var{x0}, \var{y0}, \var{level}) refines the
element containing the point $(x0,y0)$; in 3D the point has a third
coordinate \var{z0}.
\item \method{refineRegion}(\var{x0}, \var{y0}, \var{x1}, \var{y1},
\var{level}) refines the elements whose lower left corner lies in the
rectangle between $(x0,y0)$ and $(x1,y1)$; in 3D the box also takes
\var{z0} and \var{z1}.
\item \method{refineCircle}(\var{x0}, \var{y0}, \var{r}, \var{level})
(2D only) refines the elements which overlap the disc of radius \var{r}
around $(x0,y0)$.
\item \method{refineSphere}(\var{x0}, \var{y0}, \var{z0}, \var{r},
\var{level}) (3D only) refines the elements which overlap the ball of radius
\var{r} around $(x0,y0,z0)$.
\item \method{refineBorder}(\var{border}, \var{dx}, \var{level}) refines the
elements which reach into the layer of depth \var{dx} along a face of the
domain. In 2D, \var{border} is one of \var{"top"} and \var{"bottom"} (or
\var{"north"} and \var{"south"}), the faces normal to $x_1$, and
\var{"left"} and \var{"right"} (or \var{"west"} and \var{"east"}), the
faces normal to $x_0$. In 3D, \var{"top"} and \var{"bottom"} are the faces
normal to $x_2$, \var{"north"} and \var{"south"} (or \var{"back"} and
\var{"front"}) those normal to $x_1$, and \var{"east"} and \var{"west"} (or
\var{"right"} and \var{"left"}) those normal to $x_0$.
\item \method{refineMask}(\var{tag}, \var{level}) refines the elements in
which the mask named \var{tag} is positive, see
Section~\ref{sec:oxley:mask}.
\item \method{print}() lists the queued operations.
\item \method{apply}(\var{domain}, \var{**masks}) returns a new domain:
\var{domain} refined by the queued operations, with \var{masks} giving the
mask for each tag. \var{domain} is not modified.
\end{itemize}

\paragraph{toFinley()} is a method of an \oxley domain. It returns a \finley
domain describing the same mesh, with each element split into triangles (2D)
or tetrahedra (3D). \finley has to be imported before it is called.

\paragraph{toFinleyData(source, target)}
copies \var{source}, a \Data object on the \class{ContinuousFunction} of an
\oxley domain, onto \var{target}, the \finley domain \method{toFinley}
returned for it.

\paragraph{fromFinleyData(source, target)}
copies \var{source}, a \Data object on the \class{ContinuousFunction} of a
\finley domain returned by \method{toFinley}, back onto \var{target}, the
\oxley domain it was created from.
